\documentclass[11pt]{article}

% ── Packages ──────────────────────────────────────────────────────────────────
\usepackage[margin=1in]{geometry}
\usepackage{amsmath, amssymb, amsthm}
\usepackage{mathtools}
\usepackage{xcolor}
\usepackage{mdframed}
\usepackage{enumitem}
\usepackage{hyperref}
\usepackage{pagecolor}
\usepackage{microtype}

% ── Page colour (cream) ───────────────────────────────────────────────────────
\definecolor{cream}{HTML}{FFFDF5}
\definecolor{defblue}{HTML}{1A4A7A}
\definecolor{thmgreen}{HTML}{1A5C3A}
\definecolor{warnAmber}{HTML}{7A4A00}
\definecolor{dangerRed}{HTML}{7A1A1A}
\definecolor{boxbluebg}{HTML}{EBF3FB}
\definecolor{boxgreenbg}{HTML}{EBF7EF}
\definecolor{boxamberbg}{HTML}{FDF5E6}
\definecolor{boxredbg}{HTML}{FBEBEB}
\pagecolor{cream}

% ── Theorem environments ───────────────────────────────────────────────────────
\newtheoremstyle{defstyle}{8pt}{8pt}{\normalfont}{0pt}{\bfseries\color{defblue}}{.}{0.5em}{}
\newtheoremstyle{thmstyle}{8pt}{8pt}{\normalfont}{0pt}{\bfseries\color{thmgreen}}{.}{0.5em}{}
\newtheoremstyle{notstyle}{8pt}{8pt}{\normalfont}{0pt}{\bfseries\color{warnAmber}}{.}{0.5em}{}
\newtheoremstyle{cautionstyle}{8pt}{8pt}{\normalfont}{0pt}{\bfseries\color{dangerRed}}{.}{0.5em}{}

\theoremstyle{defstyle}
\newtheorem{definition}{Definition}[section]

\theoremstyle{thmstyle}
\newtheorem{theorem}{Theorem}[section]
\newtheorem{corollary}[theorem]{Corollary}

\theoremstyle{notstyle}
\newtheorem*{remark}{Remark}
\newtheorem*{example}{Example}

\theoremstyle{cautionstyle}
\newtheorem*{caution}{Caution}

% ── Coloured framed boxes ──────────────────────────────────────────────────────
\newmdenv[
  backgroundcolor=boxbluebg,
  linecolor=defblue,
  linewidth=2pt,
  topline=false, rightline=false, bottomline=false,
  innerleftmargin=10pt, innerrightmargin=10pt,
  innertopmargin=6pt, innerbottommargin=6pt,
  skipabove=8pt, skipbelow=8pt
]{defbox}

\newmdenv[
  backgroundcolor=boxgreenbg,
  linecolor=thmgreen,
  linewidth=2pt,
  topline=false, rightline=false, bottomline=false,
  innerleftmargin=10pt, innerrightmargin=10pt,
  innertopmargin=6pt, innerbottommargin=6pt,
  skipabove=8pt, skipbelow=8pt
]{thmbox}

\newmdenv[
  backgroundcolor=boxamberbg,
  linecolor=warnAmber,
  linewidth=2pt,
  topline=false, rightline=false, bottomline=false,
  innerleftmargin=10pt, innerrightmargin=10pt,
  innertopmargin=6pt, innerbottommargin=6pt,
  skipabove=8pt, skipbelow=8pt
]{notebox}

\newmdenv[
  backgroundcolor=boxredbg,
  linecolor=dangerRed,
  linewidth=2pt,
  topline=false, rightline=false, bottomline=false,
  innerleftmargin=10pt, innerrightmargin=10pt,
  innertopmargin=6pt, innerbottommargin=6pt,
  skipabove=8pt, skipbelow=8pt
]{cautionbox}

% ── Shorthand macros ──────────────────────────────────────────────────────────
\newcommand{\Z}{\mathbb{Z}}
\newcommand{\N}{\mathbb{N}}
\newcommand{\Zn}{\mathbb{Z}_n}
\newcommand{\modn}{\pmod{n}}
\newcommand{\congn}{\equiv}

% ── Document ──────────────────────────────────────────────────────────────────
\begin{document}

\begin{center}
  {\Large\bfseries Lecture 24 --- Modular Congruence}\\[4pt]
  {\normalsize CS 173: Discrete Structures \quad$\cdot$\quad Number Theory}\\[2pt]
  {\small\color{gray} University of Illinois at Urbana-Champaign}
\end{center}

\vspace{1em}
\hrule
\vspace{1.5em}

% ══════════════════════════════════════════════════════════════════════════════
\section{Definition of Congruence Modulo \texorpdfstring{$n$}{n}}

\begin{defbox}
\begin{definition}[Congruence modulo $n$]
Given $x, y \in \Z$ and $n \in \N^+$, we say $x$ is \emph{congruent to $y$ modulo $n$}, written
\[
  x \equiv y \pmod{n},
\]
if and only if $n \mid (x - y)$.
\end{definition}
\end{defbox}

\begin{notebox}
\begin{remark}
The symbol $\equiv$ here denotes an \emph{equivalence relation} on $\Z$.  Think of it as a structured version of $=$: rather than asking whether two integers are identical, we ask whether they leave the same remainder when divided by $n$.  The biconditional $p \equiv q$ uses $=$ as the underlying sign of equality at the level of equivalence classes.
\end{remark}
\end{notebox}

% ══════════════════════════════════════════════════════════════════════════════
\section{Congruence is an Equivalence Relation}

\begin{thmbox}
\begin{theorem}[Reflexivity]
\[
  \forall x \in \Z,\quad x \equiv x \pmod{n}.
\]
\begin{proof}
$x - x = 0$ and $n \mid 0$ for all $n \in \N^+$.
\end{proof}
\end{theorem}
\end{thmbox}

\begin{thmbox}
\begin{theorem}[Symmetry]
\[
  \forall x, y \in \Z,\quad x \equiv y \pmod{n} \;\Rightarrow\; y \equiv x \pmod{n}.
\]
\begin{proof}
If $n \mid (x-y)$, then $n \mid (-(x-y)) = (y-x)$.
\end{proof}
\end{theorem}
\end{thmbox}

\begin{thmbox}
\begin{theorem}[Transitivity]
\[
  \forall x,y,z \in \Z,\quad \bigl(x \equiv y \pmod{n} \;\wedge\; y \equiv z \pmod{n}\bigr)
  \;\Rightarrow\; x \equiv z \pmod{n}.
\]
\begin{proof}
$n \mid (x-y)$ and $n \mid (y-z)$ imply $n \mid \bigl((x-y)+(y-z)\bigr) = x-z$.
\end{proof}
\end{theorem}
\end{thmbox}

% ══════════════════════════════════════════════════════════════════════════════
\section{Congruence and Addition}

\begin{thmbox}
\begin{theorem}[Addition compatibility]
\[
  (\forall n \in \N^+)\;(\forall x,y \in \Z)\;(\forall z \in \Z)\qquad
  x \equiv y \pmod{n} \;\Rightarrow\; x+z \equiv y+z \pmod{n}.
\]
\end{theorem}
\end{thmbox}

\begin{thmbox}
\begin{theorem}[Additive cancellation]
Addition cancellation holds unconditionally:
\[
  x + z \equiv y + z \pmod{n} \;\Rightarrow\; x \equiv y \pmod{n}.
\]
\end{theorem}
\end{thmbox}

% ══════════════════════════════════════════════════════════════════════════════
\section{Congruence and Multiplication}

\begin{thmbox}
\begin{theorem}[Multiplication compatibility]
\[
  (\forall n \in \N^+)\;(\forall x,y \in \Z)\;(\forall z \in \Z)\qquad
  x \equiv y \pmod{n} \;\Rightarrow\; x \cdot z \equiv y \cdot z \pmod{n}.
\]
\end{theorem}
\end{thmbox}

\begin{cautionbox}
\begin{caution}[Multiplicative cancellation is \emph{not} always valid]
The converse requires an extra hypothesis.  In general,
\[
  x \cdot z \equiv y \cdot z \pmod{n}
  \;\;\text{and}\;\; z \not\equiv 0 \pmod{n}
  \;\;\not\Rightarrow\;\; x \equiv y \pmod{n}.
\]
Cancellation \emph{is} valid when $\gcd(z, n) = 1$:
\[
  x \cdot z \equiv y \cdot z \pmod{n}
  \;\;\text{and}\;\; \gcd(z,n)=1
  \;\;\Rightarrow\;\; x \equiv y \pmod{n}.
\]
\end{caution}
\end{cautionbox}

% ══════════════════════════════════════════════════════════════════════════════
\section{The Set \texorpdfstring{$\mathbb{Z}_n$}{Zn} and Modular Arithmetic}

\subsection{Definition of \texorpdfstring{$\mathbb{Z}_n$}{Zn}}

\begin{defbox}
\begin{definition}[Integers modulo $n$]
For $n \in \N^+$, define
\[
  \mathbb{Z}_n \;=\; \{0, 1, 2, \ldots, n-1\}.
\]
This is the set of \emph{residues} modulo $n$.  Every integer $x \in \Z$ is congruent to exactly one element of $\mathbb{Z}_n$.
\end{definition}
\end{defbox}

\subsection{Modular Addition and Multiplication on \texorpdfstring{$\mathbb{Z}_n$}{Zn}}

\begin{defbox}
\begin{definition}[Operations on $\mathbb{Z}_n$]
For $a, b \in \mathbb{Z}_n$, define:
\begin{align*}
  a \oplus b &\;=\; (a + b) \bmod n, \\
  a \otimes b &\;=\; (a \cdot b) \bmod n.
\end{align*}
Results are always reduced to lie in $\{0, 1, \ldots, n-1\}$.
\end{definition}
\end{defbox}

\begin{example}
Working in $\mathbb{Z}_6 = \{0,1,2,3,4,5\}$:
\begin{align*}
  1 \cdot 2 &= 2, \\
  2 \cdot 2 &= 4, \\
  4 \cdot 4 &\equiv -2 \equiv 4 \pmod{6}, \\
  2 + 1 &= 3, \\
  2 + 4 &\equiv 0 \pmod{6}, \\
  3 + 3 &\equiv 0 \pmod{6}, \\
  1 + 5 &\equiv 0 \pmod{6}, \\
  2 + 5 &\equiv 1 \pmod{6}, \\
  3 + 0 &\equiv 3 \pmod{6}.
\end{align*}
\end{example}

\subsection{Algebraic Properties of \texorpdfstring{$(\mathbb{Z}_n, +, \cdot)$}{(Zn, +, .)}}

$\mathbb{Z}_n$ equipped with modular addition and multiplication satisfies the following properties.

\begin{thmbox}
\begin{theorem}[Commutativity]
For all $a, b \in \mathbb{Z}_n$:
\[
  a \oplus b = b \oplus a, \qquad a \otimes b = b \otimes a.
\]
\end{theorem}
\end{thmbox}

\begin{thmbox}
\begin{theorem}[Associativity]
For all $a, b, c \in \mathbb{Z}_n$:
\[
  (a \oplus b) \oplus c = a \oplus (b \oplus c), \qquad
  (a \otimes b) \otimes c = a \otimes (b \otimes c).
\]
\end{theorem}
\end{thmbox}

\begin{thmbox}
\begin{theorem}[Distributivity]
For all $a, b, c \in \mathbb{Z}_n$:
\[
  a \otimes (b \oplus c) = (a \otimes b) \oplus (a \otimes c).
\]
\end{theorem}
\end{thmbox}

\begin{thmbox}
\begin{theorem}[Additive identity]
The element $0 \in \mathbb{Z}_n$ is the additive identity:
\[
  \forall a \in \mathbb{Z}_n, \quad a \oplus 0 = a.
\]
\end{theorem}
\end{thmbox}

\begin{thmbox}
\begin{theorem}[Additive inverse]
Every element of $\mathbb{Z}_n$ has an additive inverse.  For $a \in \mathbb{Z}_n$,
\[
  -a \;\equiv\; n - a \pmod{n}.
\]
That is, $a \oplus (n-a) \equiv 0 \pmod{n}$.
\end{theorem}
\end{thmbox}

\begin{notebox}
\begin{remark}
The presence of additive inverses is precisely what the board examples illustrate:
\[
  1+5\equiv 0,\quad 2+4\equiv 0,\quad 3+3\equiv 0 \pmod{6}.
\]
Each pair $(a,\, 6-a)$ sums to $0$ in $\mathbb{Z}_6$.
\end{remark}
\end{notebox}

\begin{thmbox}
\begin{theorem}[Multiplicative identity]
The element $1 \in \mathbb{Z}_n$ is the multiplicative identity:
\[
  \forall a \in \mathbb{Z}_n, \quad a \otimes 1 = a.
\]
\end{theorem}
\end{thmbox}

\begin{cautionbox}
\begin{caution}[Multiplicative inverses need not exist]
Unlike additive inverses, a multiplicative inverse of $a$ in $\mathbb{Z}_n$ exists if and only if $\gcd(a,n) = 1$.  When $n$ is prime, every nonzero element has a multiplicative inverse and $\mathbb{Z}_n$ forms a \emph{field}.
\end{caution}
\end{cautionbox}

% ══════════════════════════════════════════════════════════════════════════════
\section{Summary}

\begin{enumerate}[leftmargin=*, label=\arabic*.]
  \item $x \equiv y \pmod{n} \iff n \mid (x-y)$.
  \item $\equiv \pmod{n}$ is an equivalence relation (reflexive, symmetric, transitive).
  \item Addition is fully compatible: cancellation holds in both directions.
  \item Multiplication is compatible in one direction.  Cancellation requires $\gcd(z,n)=1$.
  \item $\mathbb{Z}_n = \{0,1,\ldots,n-1\}$ is closed under modular $+$ and $\cdot$.
  \item $(\mathbb{Z}_n, +)$ is a commutative group: it is associative, commutative, has an identity ($0$), and every element has an additive inverse ($n-a$).
  \item $(\mathbb{Z}_n, +, \cdot)$ satisfies distributivity, making it a commutative ring.
\end{enumerate}

\begin{notebox}
\begin{remark}[Looking ahead]
These ring-theoretic properties underpin modular arithmetic.  We will use them to solve linear congruences $ax \equiv b \pmod{n}$ via the extended Euclidean algorithm, and to characterise when solutions exist and how many there are.
\end{remark}
\end{notebox}

\end{document}